\Needspace{14\baselineskip}
\Article{9}
In derselben Abhandlung hat der ausgezeichnete Euler auf S. 263 versucht, unseren Satz noch auf einem anderen Weg zu bestätigen, dessen wesentlicher Inhalt Folgendes ist. Für eine vorgelegte Gleichung $x^n+Ax^{n-1}+Bx^{n-2}\ \text{usw.}=0$ konnte bisher zwar kein analytischer Ausdruck gefunden werden, der ihre Wurzeln darstellt, wenn der Exponent $n>4$ ist; dennoch scheint sicher zu sein (wie Euler behauptet), dass er nichts anderes enthalten kann als Rechenoperationen und Wurzelziehungen, die umso verwickelter miteinander verbunden sind, je größer $n$ ist. Gibt man dies zu, so zeigt Euler sehr gut, dass sich der Wert der Formel, wie verwickelt die Wurzelzeichen auch miteinander verbunden sein mögen, stets durch die Form $M+N\BareRoot{-1}$ darstellen lässt, wobei $M,N$ reelle Größen sind.

Gegen diese Schlussweise lässt sich einwenden, dass nach den Bemühungen so vieler großer Mathematiker nur sehr geringe Hoffnung bleibt, jemals zur allgemeinen Auflösung algebraischer Gleichungen zu gelangen, sodass es immer wahrscheinlicher wird, dass eine solche Auflösung überhaupt unmöglich und widersprüchlich ist. Dies darf umso weniger paradox erscheinen, \textit{als das, was man gewöhnlich die Auflösung einer Gleichung nennt, eigentlich nichts anderes ist als ihre Zurückführung auf reine Gleichungen}. Denn die Auflösung reiner Gleichungen wird dadurch nicht gelehrt, sondern vorausgesetzt; und wenn man eine Wurzel der Gleichung $x^m=H$ durch $\BareRoot[m]{H}$ ausdrückt, hat man sie keineswegs gelöst und nicht mehr getan, als wenn man zur Bezeichnung einer Wurzel der Gleichung $x^n+Ax^{n-1}+\text{usw.}=0$ irgendein Zeichen ersänne und die Wurzel diesem gleichsetzte. Es ist wahr, dass reine Gleichungen wegen der Leichtigkeit, ihre Wurzeln näherungsweise zu finden, und wegen des eleganten Zusammenhangs, den alle ihre Wurzeln untereinander haben, alle übrigen weit übertreffen. Daher ist keineswegs zu tadeln, dass die Analysten ihre Wurzeln durch ein besonderes Zeichen bezeichnet haben. Daraus aber, dass sie dieses Zeichen ebenso wie die Rechenzeichen für Addition, Subtraktion, Multiplikation, Division und Potenzierung unter dem Namen \textit{analytische Ausdrücke} zusammengefasst haben, folgt keineswegs, dass sich eine Wurzel jeder Gleichung durch sie darstellen lässt. Oder, um die Worte beiseitezulassen: Ohne hinreichenden Grund wird vorausgesetzt, dass sich die Auflösung jeder Gleichung auf die Auflösung reiner Gleichungen zurückführen lässt. Vielleicht wäre es nicht so schwierig, diese Unmöglichkeit schon für den fünften Grad mit aller Strenge zu beweisen; meine Untersuchungen darüber werde ich an anderer Stelle ausführlicher vorlegen.
\SourceLeaf{023}
Hier genügt es, dass die allgemeine Lösbarkeit der Gleichungen, in jenem Sinn verstanden, noch sehr zweifelhaft ist und daher ein Beweis, dessen ganze Beweiskraft von jener Voraussetzung abhängt, beim gegenwärtigen Stand der Sache kein Gewicht hat.
